% part: intuitionistic-logic
% chapter: semantics
% section: introduction

\documentclass[../../../include/open-logic-section]{subfiles}

\begin{document}

\olfileid{int}{sem}{int}

\olsection{పరిచయం}

అర్థవిజ్ఞానం లేకుండా ఏ తర్కాన్నీ సంతృప్తికరంగా
వివరించలేం; అంతఃప్రజ్ఞావాద తర్కమూ మినహాయింపు
కాదు. సాంప్రదాయిక తర్కానికి
\tetoken{విలువకేటాయింపుల}{!!{valuation}s} ఆధారిత
అర్థవిజ్ఞానం ప్రామాణికమైనది. అయితే
అంతఃప్రజ్ఞావాద తర్కానికి పరస్పరం పోటీపడే
అనేక అర్థవిజ్ఞానాలు ఉన్నాయి. సంధాయకాల అర్థాలకు
అంతఃప్రజ్ఞావాద దృక్కోణం నుంచి ఆమోదయోగ్యమైన
వివరణను ఇవ్వడం అనే అర్థంలో వాటిలో ఏదీ
పూర్తిగా సంతృప్తికరం కాదు.

మోడల్ తర్కాల అర్థవిజ్ఞానాన్ని పోలిన
\tetoken{సంబంధ నమూనాల}{!!{relational model}s}
ఆధారిత అర్థవిజ్ఞానం బహుశా అత్యంత ప్రాచుర్యంలో
ఉంది. ఇందులో \tetoken{ప్రతిజ్ఞావాక్య చరాలను}{!!{propositional variable}s}
లోకాలకు కేటాయిస్తారు; ఆ లోకాలు ప్రాప్యత
సంబంధంతో అనుసంధానమై ఉంటాయి. ఆ సంబంధం
ఎల్లప్పుడూ పాక్షిక క్రమం; అంటే అది స్వావర్తన,
ప్రతిసౌష్ఠవ, సంక్రామక ధర్మాలు కలిగి ఉంటుంది.

అంతఃప్రజ్ఞాత్మకంగా ఈ లోకాలను జ్ఞానస్థితులుగా
లేదా ``సాక్ష్య స్థితులుగా'' చూడవచ్చు. స్థితి~$w'$
అనేది~$w$ నుంచి ప్రాప్యమైనది అయితే, అప్పుడే
మనకు తెలిసినంతవరకు~$w'$ ఒక సాధ్యమైన (భవిష్యత్)
జ్ఞానస్థితి; అంటే~$w$ వద్ద తెలిసినదానికి
అనుకూలమైనది. ఒక ప్రతిపాదన తెలిసిన తరువాత
అది తిరిగి తెలియనిదిగా మారదు:~$w$ వద్ద $!A$
తెలిసి,~$Rww'$ అయితే,~$w'$ వద్ద కూడా $!A$
తెలుస్తుంది. కాబట్టి ``జ్ఞానం'' ప్రాప్యత సంబంధానికి
సంబంధించి ఏకదిశగా పెరుగుతుంది.

జ్ఞానసంబంధ తర్కంలోలాగా ``$!A$ తెలిసింది'' అంటే
``అన్ని జ్ఞానపర ప్రత్యామ్నాయాల్లో సత్యం'' అని
నిర్వచిస్తే, అన్ని ప్రత్యామ్నాయాల్లో $!A$,~$!B$
రెండూ తెలిసినప్పుడు $!A \land !B$ కూడా~$w$
వద్ద తెలుస్తుంది. అయితే జ్ఞానం ఏకదిశగా
పెరగడం, $R$ స్వావర్తనంగా ఉండడం వల్ల,
$!A$, $!B$ రెండూ~$w$ వద్ద తెలిసినప్పుడు,
అప్పుడే $!A \land
!B$ అనేది~$w$ వద్ద తెలుస్తుంది. ఇదే కారణంతో
వాటిలో కనీసం ఒకటి తెలిసినప్పుడు,
అప్పుడే $!A \lor !B$ అనేది $w$ వద్ద
తెలుస్తుంది. కాబట్టి $\land$, $\lor$లకు
సత్య షరతులు సాంప్రదాయిక తర్కంలోని వాటితో
సరిపోతాయి.

కానీ సోపాధికానికి సత్య షరతులు సాంప్రదాయిక
తర్కానికి భిన్నంగా ఉంటాయి. $Rww'$ అయ్యే ఏ~$w'$
వద్దా $!B$ తెలియకుండా $!A$ తెలిసే పరిస్థితి
లేనప్పుడు, అప్పుడే $!A \lif !B$ అనేది~$w$
వద్ద తెలుస్తుంది. ఇది~$w$ వద్ద $!A$ తెలియదు
లేదా $!B$~తెలుసు అనే షరతుతో సమానం కాదు.
ఎందుకంటే~$w$ వద్ద $!A$, $!B$ రెండూ
తెలియకపోయినా, $Rww'$ అయ్యే భవిష్యత్
జ్ఞానస్థితి~$w'$ ఉండవచ్చు; ఆ~$w'$ వద్ద
$!B$ తెలియకుండానే $!A$ తెలిసే అవకాశం ఉంది.

$!A$ తెలిసే సాధ్యమైన భవిష్యత్ జ్ఞానస్థితి
ఏదీ లేకపోతేనే $\lnot !A$ మనకు తెలుస్తుంది.
ఇక్కడ భావన ఏమిటంటే, $!A$ తెలిసే అవకాశం
ఉంటే, ఏదో సాధ్యమైన భవిష్యత్ జ్ఞానస్థితిలో
$!A$ తెలిసిపోతుంది. మనకు $\lfalse$ తెలియడం
అసాధ్యం కాబట్టి, ఆ భవిష్యత్ స్థితిలో $!A$
తెలిసినా~$\lfalse$ తెలియదు.

ఈ వ్యాఖ్యానంలో బహిష్కృత మధ్యమ సూత్రం
నిలవదు. ఎందుకంటే ఇంకా తెలియని, కానీ
భవిష్యత్తులో తెలియవచ్చే కొన్ని $!A$లు ఉంటాయి.
అటువంటి \tetoken{ఒక సూత్రం}{!!a{formula}}~$!A$
విషయంలో $!A$,~$\lnot !A$ రెండూ తెలియవు;
కాబట్టి $!A \lor \lnot !A$ కూడా తెలియదు.
అయితే, ఉదాహరణకు, $\lnot(!A
\land \lnot !A)$ అని మనకు తెలుసు.
ఎందుకంటే $!A$, $\lnot !A$ రెండూ తెలిసే
భవిష్యత్ స్థితి సాధ్యం కాదు; $!A$ లేదా
$\lnot !A$ ప్రస్తుతం మనకు తెలుసా లేదా
అనేదానికి సంబంధం లేకుండానే ఇది తెలుస్తుంది.

అంతఃప్రజ్ఞావాద తర్కానికి
\tetoken{సంబంధ నమూనాలు}{!!^{relational model}s}
మాత్రమే అందుబాటులో ఉన్న అర్థవిజ్ఞానం కాదు.
స్థలవిజ్ఞాన అర్థవిజ్ఞానం మరొకటి: ఇందులో
ప్రతిపాదనలను స్థలవిజ్ఞాన అంతరాళంలోని
వివృత సమితులుగా, సంధాయకాలను ఆ సమితులపై
చర్యలుగా అర్థం చేసుకుంటాం (ఉదాహరణకు
$\land$ ఛేదనానికి అనురూపం).

\end{document}
